\documentclass[11pt]{article}
\usepackage[a4paper,margin=25mm]{geometry}
\usepackage[T1]{fontenc}
\usepackage{lmodern,microtype,amsmath,amssymb,amsthm,mathtools}
\usepackage[hidelinks]{hyperref}
\hypersetup{pdftitle={Local Thermal Faithfulness from Bounded Coupling},pdfauthor={ASTRA research draft},pdfsubject={Self-contained spectral-compression bounds; priority and independent validation unestablished}}
\usepackage{enumitem}
\usepackage{fancyhdr}
\pagestyle{fancy}\fancyhf{}\fancyhead[L]{\small ASTRA THEORETICAL}\fancyhead[R]{\small Local thermal faithfulness}\fancyfoot[C]{\thepage}
\setlength{\headheight}{14pt}
\setlength{\parskip}{4pt}
\setlength{\parindent}{0pt}
\newtheorem{theorem}{Theorem}[section]
\newtheorem{lemma}[theorem]{Lemma}
\newtheorem{corollary}[theorem]{Corollary}
\newtheorem{proposition}[theorem]{Proposition}
\theoremstyle{remark}\newtheorem{remark}[theorem]{Remark}
\newcommand{\Tr}{\operatorname{Tr}}
\newcommand{\op}{\mathrm{op}}
\newcommand{\HS}{\mathrm{F}}
\newcommand{\KMS}{\mathrm{KMS}}
\newcommand{\BKM}{\mathrm{BKM}}
\newcommand{\rank}{\operatorname{rank}}
\newcommand{\supp}{\operatorname{supp}}
\newcommand{\id}{\mathbf 1}
\newcommand{\norm}[1]{\left\lVert #1\right\rVert}
\newcommand{\abs}[1]{\left\lvert #1\right\rvert}
\newcommand{\tauS}{\tau_S}
\title{Local Thermal Faithfulness from Bounded Coupling\\[4pt]\large Spectral compression, finite-band bounds, and sharp obstructions}
\author{ASTRA research draft}
\date{12 October 2026}
\begin{document}
\maketitle
\begin{abstract}
Let a finite subsystem $S$ be coupled to an arbitrary finite bath by a bounded interaction, including all subsystem-only terms. A spectral-subspace argument gives explicit lower bounds on the Kubo--Martin--Schwinger quadratic form of every local operator. The constants depend on the subsystem dimension and the interaction norm, but not on bath dimension, bath energy scale, geometry, or a spectral gap. A second theorem retains a nonzero lower bound within an energy-difference band of width seven times the interaction norm. Consequences include uniform conditioning of local-field response when the bath and coupling are fixed. Exact examples show that exponential low-temperature suppression and a subsystem-dimension dependence cannot generally be removed. The supplied mathematics is self-contained; external correctness review and comprehensive priority verification have not been completed. No full many-body learning, rapid thermalization, or total-resource technological advantage is established.
\end{abstract}

\section{Question, scope, and prior mathematics}
The investigated obstruction is quantitative loss of local thermal information: can a large, complicated environment make a nonzero local operator arbitrarily invisible to thermal quadratic forms while its coupling to the measured subsystem remains bounded? The stronger eventual target would control the covariance of an extensive sum of local Hamiltonian terms, uniformly in system size. The latter target is not proved here.

Polynomial-time Hamiltonian learning at fixed, arbitrary inverse temperature is already available \cite{BLMT}; local learning with near-optimal dependence on size and accuracy on lattices is also available \cite{CAN}. Thus, merely proving finite-temperature identifiability would not establish a new broad learning capability. The precise comparator used here is the explicit local KMS-faithfulness estimate in \cite[Lemma III.6]{CAN}. In their normalization, with interaction degree $d_g$, it states
\begin{equation}\label{eq:prior}
 \norm{B}_{\tau}\leq
 e^{80\beta\abs{\supp B}+16d_g\beta\log(2d_g\beta)}\norm{B}_{\rho,\KMS},
 \qquad \beta\geq (4d_g)^{-1}.
\end{equation}
Bergamaschi's October 2026 stationarity-test manuscript still invokes this estimate in Lemma 3.6 \cite{Ber}. That manuscript also establishes structure learning and learning from specified metastable states, so those are not claimed as newly available here. The comparison below concerns the faithfulness inequality, not the full complexities of either algorithm.

Repository checks used these revisions:\par
{\small\begin{tabular}{@{}ll@{}}
Astra & \texttt{3a9880bf4b1913a1b9e522f15414092e57c66eec}\\
OpenAI/math & \texttt{fd4aeeb2ee4fc729c18d98444fed42fd0529eeeb}
\end{tabular}}\par
The checked material was Astra's entry point, lexical routing, and scoped dossier N578, and OpenAI/math's README and targeted Hamiltonian/KMS searches. No result from either archive is used as a premise. Comparison with the broader operator-algebra and bounded-perturbation literature remains necessary before claiming priority.

\section{Definitions and main bound}
Throughout, Hilbert spaces are finite-dimensional, all traces without a subscript are ordinary unnormalized matrix traces, and $\hbar=1$. Let
\begin{equation}\label{eq:setup}
 \mathcal H=\mathcal H_S\otimes\mathcal H_B,\qquad
 d=\dim\mathcal H_S,\qquad H=H_0+V,\qquad H_0=I_S\otimes K,
\end{equation}
where $K=K^\dagger$, $V=V^\dagger$, and $\norm V_\op\leq v$. For $\beta>0$, write
\[
 Z=\Tr e^{-\beta H},\quad Z_B=\Tr_B e^{-\beta K},\quad
 \rho=Z^{-1}e^{-\beta H},\quad \tauS(Y)=\frac{1}{d}\Tr_S Y.
\]
For a local, not necessarily Hermitian operator $x$ put $X=x\otimes I_B$. Define
\begin{equation}
 \mathcal K_\rho(X)=\norm X_{\rho,\KMS}^2
 =\Tr(\rho^{1/2}X^\dagger\rho^{1/2}X).
\end{equation}
If $H|i\rangle=E_i|i\rangle$ and $p_i=e^{-\beta E_i}/Z$, then
\begin{equation}\label{eq:kms-spectral}
 \mathcal K_\rho(X)=\sum_{i,j}\sqrt{p_i p_j}\abs{X_{ij}}^2.
\end{equation}
The energy-band restriction is
\begin{equation}
 \mathcal K_{\rho,\Omega}(X)=
 \sum_{\abs{E_i-E_j}\leq\Omega}\sqrt{p_i p_j}\abs{X_{ij}}^2.
\end{equation}
It is basis-independent within degenerate eigenspaces.

\begin{theorem}[Local thermal faithfulness]\label{thm:main}
Under \eqref{eq:setup}, for every $\lambda>\sqrt3$,
\begin{equation}\label{eq:main-lambda}
 \mathcal K_\rho(X)\geq
 \frac{(1-3/\lambda^2)^2}{d}
 e^{-(\lambda+1)\beta v}\norm x_\op^2.
\end{equation}
In particular,
\begin{equation}\label{eq:main}
 \boxed{\mathcal K_\rho(X)\geq
 \frac{4}{9d}e^{-4\beta v}\norm x_\op^2
 \geq\frac{4}{9d}e^{-4\beta v}\tauS(x^\dagger x).}
\end{equation}
There is no restriction on $\dim\mathcal H_B$ or $\norm K_\op$ and no locality, commutativity, gap, or high-temperature assumption on the bath.
\end{theorem}

The coupling bound includes subsystem-only terms: an arbitrarily large local field cannot be left out of $v$. When $v=0$, $\rho=I_S/d\otimes e^{-\beta K}/Z_B$ and $\mathcal K_\rho(X)=\tauS(x^\dagger x)$ exactly.

\subsection{Spectral transfer}
\begin{lemma}[One-sided spectral-subspace transfer]\label{lem:transfer}
Let $A,B$ be Hermitian with $\norm{A-B}_\op\leq v$. For real $E$ and $g>0$, set
\[
 P_0=\id_{(-\infty,E]}(B),\qquad Q=\id_{(E+g,\infty)}(A).
\]
Then $\norm{QP_0}_\op\leq v/g$.
\end{lemma}
\begin{proof}
View $T=QP_0$ as a map from $\operatorname{Ran}P_0$ to $\operatorname{Ran}Q$. The compressed operators satisfy
\[
 A_QT-TB_{P_0}=Q(A-B)P_0=:R.
\]
Their spectra are separated by at least $g$. The Sylvester equation has the integral solution
\[
 T=\int_0^\infty e^{-t(A_Q-EI)}R e^{t(B_{P_0}-EI)}\,dt.
\]
This identity follows by differentiation and integration at the endpoints; both semigroups have the required exponential norm bounds. Consequently,
\[
 \norm T_\op\leq\int_0^\infty e^{-tg}v\,dt=v/g.
\]
The zero-rank cases are immediate.
\end{proof}

\begin{lemma}[Quadratic compression loss]\label{lem:compression}
Let $P_0,P$ be orthogonal projections and suppose $X$ commutes with $P_0$. Put $Q=I-P$ and $\delta=\norm{QP_0}_\op$. Then
\begin{equation}\label{eq:compression}
 \norm{P_0PXP P_0-P_0XP_0}_\op\leq3\delta^2\norm X_\op.
\end{equation}
\end{lemma}
\begin{proof}
Expand the difference as
\[
 -P_0QXP_0-P_0XQP_0+P_0QXQP_0.
\]
Since $[X,P_0]=0$, each of the first two terms has norm at most $\norm{P_0QP_0}_\op\norm X_\op$. The identity
\[
 \norm{P_0QP_0}_\op=\norm{QP_0}_\op^2=\delta^2
\]
follows from $Q^2=Q$. The third term has norm at most $\delta^2\norm X_\op$ as well. Sum the three bounds.
\end{proof}

\subsection{Thermal integration of retained singular values}
\begin{proof}[Proof of Theorem \ref{thm:main}]
Assume $v>0$; the decoupled case was evaluated above. For each $E\in\mathbb R$, let
\[
 P_0(E)=I_S\otimes\id_{(-\infty,E]}(K),\qquad
 P(E+\lambda v)=\id_{(-\infty,E+\lambda v]}(H),
\]
and put $r(E)=\rank\id_{(-\infty,E]}(K)$. Lemma \ref{lem:transfer} gives $\delta\leq1/\lambda$. Lemma \ref{lem:compression} gives compression loss at most $3\norm x_\op/\lambda^2$.

On $\operatorname{Ran}P_0(E)$, the operator $P_0XP_0$ is $r(E)$ copies of $x$. Thus its largest singular value $\norm x_\op$ occurs at least $r(E)$ times. The singular-value perturbation inequality shows that $P_0PXP P_0$ has at least $r(E)$ singular values at least $(1-3/\lambda^2)\norm x_\op$. Compression cannot increase Frobenius norm. Therefore, writing $P_H(t)=\id_{(-\infty,t]}(H)$,
\begin{equation}\label{eq:count}
 F(t):=\norm{P_H(t)XP_H(t)}_\HS^2,
 \qquad
 F(E+\lambda v)\geq r(E)(1-3/\lambda^2)^2\norm x_\op^2.
\end{equation}
In the energy eigenbasis, integration of this nonnegative counting function gives
\begin{align}
 \int_{-\infty}^{\infty}\beta e^{-\beta t}F(t)\,dt
 &=\sum_{i,j}e^{-\beta\max(E_i,E_j)}\abs{X_{ij}}^2\\
 &\leq\sum_{i,j}e^{-\beta(E_i+E_j)/2}\abs{X_{ij}}^2
 =Z\mathcal K_\rho(X).
\end{align}
Meanwhile, integrate \eqref{eq:count} and use $t=E+\lambda v$ to get
\[
 Z\mathcal K_\rho(X)\geq
 (1-3/\lambda^2)^2 e^{-\lambda\beta v}\norm x_\op^2
 \int_{-\infty}^{\infty}\beta e^{-\beta E}r(E)\,dE.
\]
The remaining integral equals $Z_B$. Finally, the min--max principle and $\norm V_\op\leq v$ imply the scalar partition-function comparison
\[
 Z\leq e^{\beta v}\Tr e^{-\beta H_0}=d e^{\beta v}Z_B.
\]
Divide by $Z$ to obtain \eqref{eq:main-lambda}. Choosing $\lambda=3$ proves \eqref{eq:main}; the last step uses $\tauS(x^\dagger x)\leq\norm x_\op^2$.
\end{proof}

\begin{remark}
The proof does not infer an operator-order comparison between $e^{-\beta H}$ and $e^{-\beta H_0}$. Such an inference from $-vI\preceq V\preceq vI$ would be invalid in general. Only the scalar trace comparison, justified by ordered eigenvalues, is used.
\end{remark}

\begin{corollary}[Singular-profile version]\label{cor:profile}
If $s_1(x)\geq\cdots\geq s_d(x)$ are the singular values of $x$, then for every $\lambda>0$,
\begin{equation}
 \mathcal K_\rho(X)\geq
 \frac{e^{-(\lambda+1)\beta v}}{d}
 \sum_{a=1}^d\left(s_a(x)-\frac{3}{\lambda^2}\norm x_\op\right)_+^2.
\end{equation}
\end{corollary}
\begin{proof}
Retain all singular values in \eqref{eq:count}, with their $r(E)$-fold multiplicities, instead of retaining only the largest. The rest of the integration is unchanged. For $v=0$, the stated inequality follows directly from the exact decoupled formula.
\end{proof}

\section{A finite energy band retains local fluctuations}
\begin{theorem}[Finite-band faithfulness]\label{thm:band}
In the same setting,
\begin{equation}\label{eq:band}
 \boxed{\mathcal K_{\rho,7v}(X)\geq
 \frac{e^{-5\beta v}}{72d}\norm x_\op^2.}
\end{equation}
In particular, the bath cannot move all local KMS fluctuation weight to arbitrarily high energy differences at fixed $d,\beta,v$.
\end{theorem}
\begin{proof}
For $v>0$, partition the bath spectrum into intervals $I_j=[jv,(j+1)v)$ and set
\[
 P_{0j}=I_S\otimes\id_{I_j}(K),\qquad
 P_j=\id_{[(j-3)v,(j+4)v]}(H),\qquad
 r_j=\rank\id_{I_j}(K).
\]
Apply the Sylvester argument separately above and below the expanded interval. Each discarded part has norm at most $1/3$ on $\operatorname{Ran}P_{0j}$. Their ranges are orthogonal, so
\[
 \norm{(I-P_j)P_{0j}}_\op^2\leq2/9.
\]
Lemma \ref{lem:compression} therefore gives compression error at most $2\norm x_\op/3$. Singular-value counting gives
\[
 \norm{P_jXP_j}_\HS^2\geq \frac{r_j}{9}\norm x_\op^2.
\]
For two energies in the $j$th expanded interval,
$e^{-\beta(E_a+E_b)/2}\geq e^{-\beta(j+4)v}$. Thus
\[
 \sum_j\sum_{a,b:\,E_a,E_b\in[(j-3)v,(j+4)v]}
 e^{-\beta(E_a+E_b)/2}\abs{X_{ab}}^2
 \geq\frac{\norm x_\op^2}{9}\sum_jr_j e^{-\beta(j+4)v}.
\]
Every energy pair in this sum has difference at most $7v$, and any pair belongs to at most eight of the closed expanded intervals. Therefore the left side is at most $8Z\mathcal K_{\rho,7v}(X)$.

For $\kappa\in I_j$, $(j+4)v\leq\kappa+4v$. Hence the right side is at least
$e^{-4\beta v}Z_B\norm x_\op^2/9$. Divide by eight and use
$Z\leq d e^{\beta v}Z_B$. For $v=0$, all local transitions have zero energy difference, and the exact decoupled formula implies the claimed bound.
\end{proof}

This is an absolute floor, not a claim that a constant fraction of the total fluctuation lies in the band. A zero-frequency contribution is allowed. Consequently, it is not a lower bound on nonzero-frequency absorption, a transport coefficient, a mixing rate, or the time needed to resolve individual spectral lines.

\section{Local-field conditioning and its operational boundary}
For a Hermitian $A$, let $A_c=A-\Tr(\rho A)I$. Define the Bogoliubov--Kubo--Mori covariance
\[
 \operatorname{Cov}^{\BKM}_\rho(A,A)=
 \int_0^1\Tr(\rho^s A_c\rho^{1-s}A_c)\,ds.
\]
In the energy eigenbasis its coefficient for a matrix entry is the logarithmic mean
$L(p_i,p_j)=\int_0^1p_i^s p_j^{1-s}\,ds$, which is at least $\sqrt{p_i p_j}$ by Jensen's inequality. Therefore
\begin{equation}\label{eq:bkm}
 \operatorname{Cov}^{\BKM}_\rho(A,A)\geq\mathcal K_\rho(A_c).
\end{equation}

\begin{corollary}[Bath-independent local response conditioning]\label{cor:inverse}
Let $F_1,\ldots,F_q$ be traceless Hermitian operators on $S$ with $\tauS(F_aF_b)=\delta_{ab}$. Fix $K$ and $V_0$, and consider
\[
 H(h)=I_S\otimes K+V_0+\sum_{a=1}^q h_a F_a\otimes I_B.
\]
Suppose $h$ lies in a convex parameter set on which
$\norm{V_0+\sum_a h_aF_a\otimes I_B}_\op\leq v$.
Let $m_a(h)=\Tr[\rho(h)(F_a\otimes I_B)]$ and
$\kappa=4e^{-4\beta v}/(9d)$. Then
\begin{align}
 \nabla^2\log Z(h)&\succeq\beta^2\kappa I_q,\label{eq:hessian}\\
 -(h-h')\cdot(m(h)-m(h'))&\geq\beta\kappa\norm{h-h'}_2^2,\label{eq:monotone}\\
 \norm{h-h'}_2&\leq\frac{9d}{4\beta}e^{4\beta v}\norm{m(h)-m(h')}_2.\label{eq:inverse}
\end{align}
In particular the local expectation map is injective on this set.
\end{corollary}
\begin{proof}
Differentiating the matrix exponential under the trace gives
$\partial_a\log Z=-\beta m_a$ and
\[
 u^\top\nabla^2\log Z\,u=
 \beta^2\operatorname{Cov}^{\BKM}_\rho(A_u,A_u),
 \quad A_u=\sum_a u_aF_a\otimes I_B.
\]
Tracial orthonormality gives $\tauS((\sum_a u_aF_a)^2)=\norm u_2^2$.
Since $A_u$ is tracially centered, subtracting its thermal mean only increases the local tracial second moment. Combine \eqref{eq:bkm} and Theorem \ref{thm:main} to prove \eqref{eq:hessian}. Integrate the Hessian along the segment from $h'$ to $h$ to obtain \eqref{eq:monotone}. Cauchy--Schwarz gives \eqref{eq:inverse}.
\end{proof}

\subsection{What a sample bound would, and would not, provide}
Assume independent exact equilibrium copies are supplied and $\norm{F_a}_\op\leq L$. Estimating each $m_a$ from an independent batch of bounded-outcome measurements gives a Euclidean moment error at most $\delta$ with failure probability at most $\alpha$ using at most
\[
 q+\frac{2q^2L^2}{\delta^2}\log\frac{2q}{\alpha}
\]
copies, by the scalar Hoeffding bound and a union bound. If a model fit has residual at most $\delta$, \eqref{eq:inverse} gives parameter error at most $2\delta/(\beta\kappa)$. Taking $\delta=\beta\kappa\epsilon/2$ yields the sufficient measurement count
\begin{equation}\label{eq:samples}
 q+\frac{81q^2L^2d^2}{2\beta^2\epsilon^2}
 e^{8\beta v}\log\frac{2q}{\alpha}.
\end{equation}
This is deliberately only a measurement and conditioning statement. Computing the forward map or minimizing an objective involving $\log Z$ has not been made efficient. If one physical equilibrium copy costs $T_{\rm eq}$, preparation alone costs the number of copies times $T_{\rm eq}$; neither the theorem nor \eqref{eq:samples} bounds that quantity. Independence, preparation bias, observable implementation, and numerical optimization must be charged separately. For Pauli observables on $k$ qubits, measurement can be local to those qubits, but arbitrary $F_a$ require their own measurement implementation.

\subsection{Unknown environments are not automatically identifiable}
The bath and $V_0$ are fixed in Corollary \ref{cor:inverse}. A simple commuting example proves that this restriction matters. With a bath qubit, take
\[
 H=hZ_S+bZ_B+JZ_SZ_B.
\]
The subsystem probabilities obey
\[
 \frac{p_+}{p_-}=e^{-2\beta h}
 \frac{\cosh\beta(b+J)}{\cosh\beta(b-J)}.
\]
Thus the local Gibbs state is maximally mixed when
\[
 h=\frac{1}{2\beta}\log
 \frac{\cosh\beta(b+J)}{\cosh\beta(b-J)},
\]
which is nonzero for suitable $b,J$. It is the same local state as an uncoupled subsystem with $h=0$. Consequently local equilibrium data alone cannot identify an unknown microscopic local field when the bath and coupling can vary simultaneously.

\section{Necessary losses: exact examples}
\subsection{Exponential low-temperature suppression}
For a single qubit with $H(h)=hZ$,
\[
 m(h)=-\tanh(\beta h),\qquad
 -m'(h)=\beta\operatorname{sech}^2(\beta h).
\]
At $h=v>0$ this is asymptotic to $4\beta e^{-2\beta v}$ as $\beta v\to\infty$. Thus a universal inverse-response guarantee polynomial in $\beta v$ is impossible even without a bath. The present bounds have the correct exponential-versus-polynomial temperature class, but their constants in the exponent and their statistical costs are not claimed optimal.

\subsection{No subsystem-dimension-free KMS comparison}
\begin{proposition}\label{prop:dimension}
At fixed $\beta=v=1$, no positive constant $c$ independent of $d$ can guarantee
\[
 \mathcal K_\rho(x\otimes I_B)\geq c\tauS(x^2)
\]
for all traceless Hermitian local $x$ and all baths in \eqref{eq:setup}. In an explicit family,
\[
 \frac{\mathcal K_\rho(x\otimes I_B)}{\tauS(x^2)}
 \leq\frac{6(4\log d+1)^4}{d}\longrightarrow0.
\]
\end{proposition}
\begin{proof}
Take even $d\geq4$, subsystem basis $|0\rangle,\ldots,|d-1\rangle$, and a bath of dimension $d-1$ with basis $|0\rangle,\ldots,|d-2\rangle$. Put
\[
 M=4\log d,\qquad
 r=\frac{2}{\sqrt{M^2+4}+M},\qquad M=r^{-1}-r,
\]
\[
 K=M\sum_{b=1}^{d-2}|b\rangle\langle b|.
\]
Let $s_j=+1$ for even $j\geq2$ and $s_j=-1$ for odd $j\geq2$, and define
\[
 x=|0\rangle\langle0|-|1\rangle\langle1|
   +r^2\sum_{j=2}^{d-1}s_j|j\rangle\langle j|.
\]
This operator is traceless and has norm one. Set $a_j=1$ when $s_j=+1$ and $a_j=0$ when $s_j=-1$. Couple only the disjoint pairs
\[
 V=\sum_{j=2}^{d-1}\bigl(|j,0\rangle\langle a_j,j-1|
             +|a_j,j-1\rangle\langle j,0|\bigr).
\]
Each block of $V$ is a Pauli $X$, so $\norm V_\op=1$.

Each interacting block of $H$ is $\left(\begin{smallmatrix}0&1\\1&M\end{smallmatrix}\right)$, with eigenvalues $-r$ and $r^{-1}$. In its energy eigenbasis the local operator $x\otimes I_B$ has low-state diagonal entry zero, high-state diagonal entry of magnitude $1-r^2$, and off-diagonal entry of magnitude $r$. The two uncoupled ground-bath states have diagonal entries $+1$ and $-1$. At bath energy $M$, there are $d-2$ uncoupled entries of magnitude one and $(d-2)^2$ of magnitude $r^2$. Direct summation gives
\begin{align}
 Z={}&2+(d-2)(e^r+e^{-1/r})+(d-2)(d-1)e^{-M},\\
 Z\mathcal K_\rho(X)={}&2+(d-2)\bigl[e^{-1/r}(1-r^2)^2+2e^{-M/2}r^2\bigr]\nonumber\\
 &+e^{-M}\bigl[(d-2)+(d-2)^2r^4\bigr].\label{eq:counter-numerator}
\end{align}
The thermal mean is zero by the sign-pair symmetry. Since $e^{-M}=d^{-4}$, $0<r<1$, and $Z\geq d$, the numerator in \eqref{eq:counter-numerator} is less than three. Hence $\mathcal K_\rho(X)\leq3/d$. On the other hand,
\[
 \tauS(x^2)=\frac{2+(d-2)r^4}{d}\geq\frac12r^4
 \geq\frac{1}{2(M+1)^4}.
\]
Here $r\geq(M+1)^{-1}$ follows from $M\geq1$ and $\sqrt{M^2+4}\leq M+2$. Combining the bounds proves the claim.
\end{proof}

The proposition establishes necessity of some dimension dependence, and matches the $d^{-1}$ power up to logarithmic factors. It is not a counterexample to a bounded-degree, fixed-locality many-body covariance theorem: its large subsystem and its couplings do not preserve those extra structural restrictions. It also concerns KMS covariance, not an asserted dimension-free impossibility for the full BKM covariance.

\section{Consequences for local spin Hamiltonians and the remaining obstruction}
Suppose $H=\sum_T h_T$ on a spin system and
\[
 g:=\max_i\sum_{T\ni i}\norm{h_T}_\op<\infty.
\]
For a region $S$ of $k$ qubits, put all terms disjoint from $S$ in $K$ and all remaining terms in $V$. Then $v\leq gk$ and $d=2^k$. Theorem \ref{thm:main} implies
\begin{equation}\label{eq:lattice}
 \boxed{\norm X_{\tau}\leq\frac32\,2^{k/2}e^{2\beta gk}\norm X_{\rho,\KMS}.}
\end{equation}
For fixed local support and strength, the logarithm of this conditioning factor is linear in $\beta$. This improves the explicit $O(\beta\log\beta)$ exponent in the inspected bound \eqref{eq:prior}, and it does not require bounded interaction degree when the sum of incident interaction norms stays bounded. None of these statements bounds the complete cost of a learning algorithm by substitution alone.

The larger target remains a volume-independent inequality of the form
\begin{equation}\label{eq:global-target}
 \operatorname{Cov}^{\BKM}_{\rho_\beta(H)}\left(\sum_{a=1}^m u_aP_a,
 \sum_{a=1}^m u_aP_a\right)
 \geq c(\beta,J,k,\Delta)\sum_{a=1}^m u_a^2
\end{equation}
for a canonical family of distinct nonidentity, bounded-locality Pauli terms and bounded interaction degree, with $c$ independent of $m$. The supplied proof does not establish \eqref{eq:global-target}. A recent bounded-degree comparison in \cite[Appendix D]{SYK} still yields a Hessian estimate with an explicit $m(1+\log(2m))$ denominator after conversion from the largest coefficient to the Euclidean coefficient norm. That observation is a comparison to that paper, not an exhaustive claim about every possible prior bound.

Why the local theorem does not suffice is precise. For an extensive sum, choosing its whole support as $S$ makes $d$ exponential and $v$ extensive. Applying the theorem separately to summands does not control cancellations among off-diagonal covariances. Positive local variances do not lower-bound the smallest eigenvalue of the full covariance matrix. The missing capability is a quantum tensorization or comparable global anti-cancellation theorem under the physical interaction assumptions. No such result is smuggled into \eqref{eq:lattice} or Corollary \ref{cor:inverse}.

\section{Verification record and status}
The accompanying \texttt{check\_bounds.py} runs 300 random finite-dimensional tests, including arbitrary non-Hermitian local operators, a range of bath energy scales, the spectral-leakage inequality, the quadratic compression inequality, the rank-profile estimate, and both thermal bounds. It also evaluates the dimension-counterexample formulas and checks three small cases by dense diagonalization. The seed is 19371. All programmed checks passed within the stated floating-point tolerances; the largest compression-inequality numerical excess was approximately $1.96\times10^{-12}$. These are falsification checks, not formal verification or independent review. Large-$d$ examples are evaluated from the proved formulas, not by constructing their enormous matrices.

The main theorem, the finite-band extension, the inverse-response corollary, and the dimension obstruction have complete supplied finite-dimensional proofs. Their original priority and the assessment that they meet an independently major foundational-completion standard remain unestablished. The global covariance target and an end-to-end technological advantage are not claimed. No new physical law or experimental hypothesis is introduced: the physical premises are an equilibrium Gibbs state and the stated Hamiltonian/coupling bounds.

\begin{thebibliography}{9}\small
\bibitem{BLMT} A. Bakshi, A. Liu, A. Moitra, and E. Tang.
\emph{Learning quantum Hamiltonians at any temperature in polynomial time}.
\href{https://arxiv.org/abs/2310.02243}{arXiv:2310.02243}, inspected version 2.
\bibitem{CAN} C.-F. Chen, A. Anshu, and Q. T. Nguyen.
\emph{Learning quantum Gibbs states locally and efficiently}.
\href{https://arxiv.org/abs/2504.02706}{arXiv:2504.02706}, version 1. In particular Lemma III.6.
\bibitem{Ber} T. Bergamaschi.
\emph{The stationarity test: a framework for learning quantum many-body systems from their thermal states}.
\href{https://arxiv.org/abs/2610.01074}{arXiv:2610.01074}, version 1; manuscript dated 2 October 2026. In particular Lemma 3.6.
\bibitem{SYK} A. Anshu, S. Arunachalam, S. Chen, and Y. Hwang.
\emph{Learning SYK Hamiltonians}.
\href{https://arxiv.org/abs/2610.02178}{arXiv:2610.02178}, version 1. In particular Appendix D.
\bibitem{ASTRA} Astra research-memory repository.
\href{https://github.com/harjassand/astra-research-knowledge}{\nolinkurl{harjassand/astra-research-knowledge}}.
Pinned revision stated in Section 1; entry point, lexical routing, and scoped dossier N578 inspected. No result imported as a premise.
\bibitem{OAI} OpenAI mathematics repository.
\href{https://github.com/openai/math}{\nolinkurl{openai/math}}.
README and targeted search inspected; search revision stated in Section 1. No result imported as a premise.
\end{thebibliography}
\end{document}
